\section{Parte 1: Diseño del Tablero Kanban en Jira}

\subsection{Estructura de Columnas y Estados del Tablero}
En cumplimiento estricto de la \textbf{Parte 1} de la guía de prácticas, se aprovisionó el \textbf{Tablero Kanban ID 137 (\textit{Tablero Kanban - Lab 06})} en Jira Software Cloud (\url{https://unsaac-team-grupois.atlassian.net/jira/software/projects/SC/boards/137}), estructurado en cinco estados:

\begin{enumerate}[leftmargin=*]
    \item \textbf{\textit{Backlog}:} Cola de requerimientos priorizados según el plan de la MDSJ.
    \item \textbf{\textit{Por hacer} (ID: 10155):} Historias refinadas listas para ser iniciadas por los desarrolladores.
    \item \textbf{\textit{En curso} (ID: 10156):} Trabajo activo de codificación en frontend/backend.
    \item \textbf{\textit{En revisión} (ID: 10158):} Control de calidad (\textit{QA}) y revisión cruzada de código (\textit{Code Review}).
    \item \textbf{\textit{Listo / Hecho} (ID: 10157):} Incrementos probados y desplegados satisfactoriamente.
\end{enumerate}

\begin{figure}[H]
    \centering
    \includegraphics[width=0.96\textwidth]{captura_jira_tablero_kanban.png}
    \caption{Tablero Kanban Activo en Jira Software Cloud (Tablero ID: 137 - SIGC-CUSCO).}
    \label{fig:tablero_kanban}
\end{figure}

\begin{table}[H]
\centering
\small
\caption{Mapeo de Estados y Configuración del Flujo en Jira.}
\label{tab:mapeo_columnas}
\begin{tabularx}{\textwidth}{P{3.2cm} P{2.5cm} P{2.8cm} Y}
\toprule
\textbf{Columna} & \textbf{ID Estado} & \textbf{Categoría} & \textbf{Definición Operativa en SIGC-CUSCO} \\
\midrule
\textit{Por hacer} & \texttt{10155} & \texttt{new} & Requerimientos priorizados de la MDSJ listos para implementar. \\
\textit{En curso} & \texttt{10156} & \texttt{indeterminate} & Desarrollo activo de controladores, migraciones y vistas React. \\
\textit{En revisión} & \texttt{10158} & \texttt{indeterminate} & Pruebas con API RENIEC y revisión de código por pares. \\
\textit{Listo} & \texttt{10157} & \texttt{done} & Módulos validados con criterios de aceptación cumplidos. \\
\bottomrule
\end{tabularx}
\end{table}
